%%%%%%%%%%%%%%%%%%%%%%%%%%%%%%%%%%%%%%%%%%%%%%%%%%%%%%%%%%%%%%%%
\section{Case-Study Objective and Validation Strategy}
\label{sec:case_study_strategy}

The case study follows from the need for a representative heterogeneous hybrid workflow described in Chapter~\ref{chap:intro}.
The controlled pendulum is simple enough to interpret but rich enough to exercise the main framework features.
It combines closed-loop control, sensor feedback, actuator dynamics, contact, and alternative pendulum models with different fidelity levels.
It also permits monolithic OpenModelica reference models for selected numerical comparisons.
These properties make it suitable for numerical verification against reference solutions and for workflow validation of the heterogeneous co-simulation framework.

The case study verifies selected co-simulation results against monolithic OpenModelica reference simulations and validates the framework workflow on a controlled pendulum system.
Performance results are reported as benchmarks.
No experimental validation is claimed.

The numerical verification follows a convergence-based strategy.
For scenarios with a monolithic reference, the co-simulation result is compared against the OpenModelica solution for decreasing macro step sizes.
The expected behavior is that the trajectory error decreases as the communication step size is reduced.
The OpenModelica models represent the same closed-loop configurations in one simulation environment.
They are numerical references and are not treated as experimental ground truth.

The workflow validation addresses the interaction between framework features.
The contact and switching scenario checks whether the framework can assemble heterogeneous components, derive execution metadata, handle contact events, route a controller reset, and switch between pendulum models during runtime.
Chapter~\ref{chap:implementation} verified these mechanisms separately.
The case study evaluates their combined use in one closed-loop workflow.

The performance benchmark addresses the computational cost of high-fidelity modeling.
The full-FEM pendulum is compared with a switched configuration that activates the FEM model only during contact-relevant phases.
The benchmark reports the trajectory deviation relative to the full-FEM run and the resulting wall-clock speedup.

The following sections keep these claim types separate.
Numerical verification is used when a monolithic OpenModelica reference is available.
Workflow validation is used when several framework features are exercised together.
Performance results are reported as benchmarks and are interpreted together with the observed numerical deviation.